%% ══════════════════════════════════════════════════════════
%%  ITEM — Shuffled Time   ·   Pocket Leaflet (A5)
%%  Compile: pdflatex item-leaflet.tex  (run twice for TOC)
%% ══════════════════════════════════════════════════════════

\documentclass[a6paper,10pt]{extarticle}

\usepackage[a6paper, top=0cm, bottom=0cm, left=0cm, right=0cm]{geometry}
\usepackage[T2A,T1]{fontenc}
\usepackage[utf8]{inputenc}
\usepackage[main=english,russian]{babel}
\usepackage{xcolor}
\usepackage{tikz}
\usepackage{graphicx}
\usepackage{parskip}
\usepackage{microtype}
\usepackage{lmodern}
\usepackage{textcomp}
\usepackage{hyperref}
\usepackage{fancyhdr}
\usepackage{eso-pic}
\usepackage{transparent}

\hypersetup{hidelinks}

\definecolor{raacream}{HTML}{F5F5DC}
\definecolor{raared}{HTML}{B29C83}
\definecolor{raadark}{HTML}{0D0D0D}
\definecolor{mygray}{HTML}{555555}

%% Image asset base path
\newcommand{\assets}{/Users/leo/Projects/Raa/web/frontend/public/assets/item}

\setlength{\parskip}{1.5pt}
\setlength{\parindent}{0pt}

%% Page style: small centred roman numeral in footer
\fancypagestyle{leaflet}{%
  \fancyhf{}%
  \renewcommand{\headrulewidth}{0pt}%
  \fancyfoot[C]{}%
}
\pagestyle{leaflet}

%% Watermark background for description pages (runs at shipout, zero text-flow impact)
\newcommand{\descpagebg}{%
  \AddToShipoutPictureBG*{%
    \put(0,0){%
      \transparent{0.80}%
      \includegraphics[width=\paperwidth,height=\paperheight,keepaspectratio=false]%
        {\assets/ITEM-background.png}%
    }%
  }%
}

%% Red section title with rule
\newcommand{\sectiontitle}[1]{%
  \vspace{4pt}%
  {\color{raared}\bfseries\small #1}\\[-4pt]%
  {\color{raared}\rule{\linewidth}{0.5pt}}%
  \vspace{2pt}%
}

%% Full-bleed image page (cover / back / posters)
%% Call after \newpage + \newgeometry{0,0,0,0}
%% Overlays a tiny white page number at bottom centre (pass {} to suppress)
\newcommand{\fullbleedimage}[2]{%
  \thispagestyle{empty}%
  \noindent%
  \begin{tikzpicture}[x=1cm,y=1cm]%
    \useasboundingbox (0,0) rectangle (10.5,14.8);%
    \node[anchor=south west, inner sep=0pt] at (0,0) {%
      \includegraphics[width=10.5cm,height=14.8cm,keepaspectratio=false]{#1}%
    };%
    \ifx&#2&\else%
      \node[anchor=south west, inner sep=0pt] at (0.4,0.3) {%
        \fontsize{6.5}{8}\selectfont\color{white}#2%
      };%
    \fi%
  \end{tikzpicture}%
}

\begin{document}
\pagenumbering{roman}

%% ══════════════════════════════════════════════════════════
%%  FRONT COVER — full-bleed poster-main.png (unnumbered)
%% ══════════════════════════════════════════════════════════
\setcounter{page}{0}
\fullbleedimage{\assets/poster-main.png}{}

%% ══════════════════════════════════════════════════════════
%%  PAGE i — Event info
%% ══════════════════════════════════════════════════════════
\newpage
\newgeometry{top=0.4cm,bottom=0.4cm,left=0.5cm,right=0.5cm}
\thispagestyle{empty}

\null\vfill

\begin{center}
  {\fontsize{28}{32}\selectfont\bfseries ITEM}\\[5pt]
  {\large\bfseries SHUFFLED TIME}\\[10pt]
  {\color{raared}\rule{4cm}{0.5pt}}\\[10pt]
  {\small Performative Week at RAA.SPACE}\\[3pt]
  {\small\textbf{5\,–\,11 July 2026}}\\[2pt]
  {\small Matīsa iela 8, Rīga}\\[10pt]
  {\color{raared}\rule{4cm}{0.5pt}}\\[10pt]
  {\small Concept and idea by\\[3pt]
   \textbf{Jana Zariņa} \& \textbf{Leonhard Horstmeyer}}\\[10pt]
  {\color{raared}\rule{4cm}{0.5pt}}\\[10pt]
  {\small With the kind support of\\[3pt]
   \textbf{Radošā Rūpnīca Veldze} and \textbf{LaiZinat.com}}\\[10pt]
  {\color{raared}\rule{4cm}{0.5pt}}\\[10pt]
  {\small
   RAA.SPACE is a free and independent cultural space.\\[3pt]
   We are grateful for your donations.\\[3pt]
   \href{https://raa.space/donations}%
     {\texttt{raa.space/donations}}}
\end{center}

\vfill

%% ══════════════════════════════════════════════════════════
%%  PAGE ii — About ITEM · About RAA · Contents
%% ══════════════════════════════════════════════════════════
\newpage
\newgeometry{top=0.4cm,bottom=0.4cm,left=0.5cm,right=0.5cm}
\thispagestyle{empty}

\sectiontitle{About ITEM}

{\small
\textbf{ITEM\,—\,Shuffled Time} is a seven-day performative week at RAA cultural
space in Rīga. Each day one artist or artist group has a 24-hour time window to
\textbf{I}nvoke, \textbf{T}emper, \textbf{E}late, \textbf{M}int the space. The
window starts and ends at 9\,h in the morning; exactly how the 24 hours are filled
is at the discretion of the artists.

Anyone is invited to pass by the space at any time. On each day there is a scheduled
group visit — see \href{https://raa.space/item}{\texttt{raa.space/item}} for times.

The seven participating artists and groups work across durational performance, body
art, experimental poetry, aromagraphy, and interactive installation. Each performance
stands independently while forming part of a shared proposition about time, space,
and presence.

{\footnotesize\color{mygray}
With thanks to Radošā Rūpnīca Veldze and LaiZinat.com.}
}

\sectiontitle{About RAA}

{\footnotesize\setlength{\parskip}{0pt}
RAA is an independent cultural space in Rīga hosting regular theatrical and
contemporary performances, exhibitions, and workshops. Located at Matīsa iela 8
within the Veldze creative quarter, the main hall covers 138\,m\textsuperscript{2}
with ten large windows facing the street. At RAA you don't have to be loud to be
heard. The name comes from \textit{Reductio ad Absurdum} — an argument that
establishes a claim by showing that its contrary leads to absurdity.

\vspace{2pt}
\textbf{raa.space}\enspace·\enspace
\href{mailto:enter@raa.space}{enter@raa.space}\enspace·\enspace
@raa\_riga
}

\sectiontitle{Contents}

{\footnotesize\setlength{\parskip}{1.5pt}
\begin{tabular*}{\textwidth}{@{}r@{\,}l@{\enspace}l@{\extracolsep{\fill}}r@{}}
2 & MON & \textbf{NOW} — Maruta Gredzena                        & \textit{iii}\\
5 & THU & \textbf{Presence in Absence Piece Three} — Lisette Ros & \textit{v}\\
3 & TUE & \textbf{N.3} — Ernests Ā                             & \textit{vii}\\
1 & SUN & \textbf{Words to Change the World} — Ilze Aulmane     & \textit{ix}\\
4 & WED & \textbf{Gluttony} — K.\,Šuksta \& S.\,Muriņš         & \textit{xi}\\
6 & FRI & \textbf{Time Scores} — Vanessa Brzić                  & \textit{xiii}\\
7 & SAT & \textbf{Laiks Telpā} — Ilze Mazpane \& Jana Zariņa   & \textit{xv}\\
\end{tabular*}
}

%% ══════════════════════════════════════════════════════════
%%  PAGES iii–iv: NOW — Maruta Gredzena  (Mon 6.07)
%% ══════════════════════════════════════════════════════════

\newpage
\newgeometry{top=0cm,bottom=0cm,left=0cm,right=0cm}
\fullbleedimage{\assets/06_07_now.png}{\thepage}

\newpage
\newgeometry{top=0.4cm,bottom=0.4cm,left=0.5cm,right=0.5cm}
\descpagebg

{\footnotesize\color{mygray}Monday, 6 July 2026}\\[1pt]
{\large\bfseries NOW}\\[2pt]
{\small\textit{Maruta Gredzena\enspace(LV)}}\\[4pt]
{\footnotesize\color{mygray}%
  \textbf{Performance:} 17:00\,–\,22:00\enspace·\enspace
  \textbf{Group visit:} 19:00}

\vspace{4pt}

{\small
In her interactive performance \textit{Now}, Maruta Gredzena explores whether
communication without words can be true. Visitors become participants, choosing and
wearing light, paper-based wearable sculptures in response to their present state.
The forms can feel like cocoons — soft, fluffy, geometric, or rock-like — creating
a temporary second skin. Each choice reshapes a shifting collective composition,
where everyone becomes both observer and artwork.

The work asks: can honesty exist without language? The forms carry the artist's
intimate visual language, yet are activated only through others. The question of
authorship and non-verbal communication arises: is the language already created in
form, or does it emerge only in the moment it is chosen and moved?

Here, form becomes language, presence becomes communication, and material becomes a
shared way of speaking.
}

\vspace{3pt}
\sectiontitle{About the artist}

{\footnotesize\setlength{\parskip}{0pt}
Maruta Gredzena is interested in what happens when nothing stays still. A life
shaped by continual relocation has shifted her practice from making objects to
creating spaces for interaction. Rooted in porcelain and ceramic sculpture and
expanded through graphic art, her practice now unfolds through installations that
use paper, material, and playfulness to explore communication beyond language.
Wandering between the Baltic coast and the Iberian Peninsula has cultivated a
migratory way of working. Drawn to places that feel temporarily safe, she treats
movement as both method and subject, transforming displacement into encounters
where form, colour, pattern, and participation become a shared language.
}

\vfill
\noindent{\footnotesize\color{mygray}\href{https://raa.space/item/now}{\texttt{raa.space/item/now}}\hfill\thepage}

%% ══════════════════════════════════════════════════════════
%%  PAGES v–vi: PRESENCE IN ABSENCE PIECE THREE — Lisette Ros  (Thu 9.07)
%% ══════════════════════════════════════════════════════════

\newpage
\newgeometry{top=0cm,bottom=0cm,left=0cm,right=0cm}
\fullbleedimage{\assets/09_07_presence_in_absence_piece_three.png}{\thepage}

\newpage
\newgeometry{top=0.4cm,bottom=0.4cm,left=0.5cm,right=0.5cm}
\descpagebg

{\footnotesize\color{mygray}Thursday, 9 July 2026}\\[1pt]
{\large\bfseries Presence in Absence Piece Three}\\[2pt]
{\small\textit{Lisette Ros\enspace(NL)}}\\[4pt]
{\footnotesize\color{mygray}%
  \textbf{Performance:} 16:00\,–\,open end\enspace·\enspace
  \textbf{Group visit:} 18:00}

\vspace{4pt}

{\small
\textit{Presence in Absence} is a research series by Lisette Ros that began in
Mexico City in 2023, deriving from her project \textit{The Absence of Metamorphosis}.
The series focuses on the perceptible act of presence in absence by examining
\textit{The Absent Body}.

In Piece Three, Ros grapples with her own presence, manipulating control and
consciousness by taking in a natural hypnotic — a substance that has formed part of
her nightly ritual since 2014, when she was diagnosed with DSPS (Delayed Sleep Phase
Syndrome). The work unfolds in real time, before the audience.
}

\vspace{3pt}
\sectiontitle{About the artist}

{\footnotesize\setlength{\parskip}{0pt}
Lisette Ros (b.\,1991, NL) is a Dutch conceptual and performance artist with
Indonesian roots. Identifying as queer, woman and humanimal, her work investigates
how identity is a construct formed through social systems, trauma, cultural
conditioning and embodied experience. Using her own body as research tool, source,
material and medium, she creates performances, video works and installations that
examine daily rituals, conventions, and the fragile borders between the personal and
the collective.

Her work has been presented internationally at the Venice Architecture Biennale
(2016), South by Southwest (Austin, 2018), EYE Filmmuseum, Van Abbemuseum, Goethe
Institute, Amsterdam Museum, MWoods Beijing, Casa Lü CDMX and Oped Space Tokyo.
In 2021 she completed a residency with Marina Abramović.

Ros is the creator of the long-term series \textit{My Self,} whose sixth chapter
\textit{My Self, the (Unfamiliar) Roots} is supported by the Mondriaan Fund and the
Dutch Embassy in Jakarta.
}

\vfill
\noindent{\footnotesize\color{mygray}%
  \href{https://raa.space/item/presenceinabsence_piecethree}%
  {\texttt{raa.space/item/presenceinabsence\_piecethree}}\hfill\thepage}

%% ══════════════════════════════════════════════════════════
%%  PAGES vii–viii: N.3 — Ernests Ā  (Tue 7.07)
%% ══════════════════════════════════════════════════════════

\newpage
\newgeometry{top=0cm,bottom=0cm,left=0cm,right=0cm}
\fullbleedimage{\assets/07_07_N_3.png}{\thepage}

\newpage
\newgeometry{top=0.4cm,bottom=0.4cm,left=0.5cm,right=0.5cm}
\descpagebg

{\footnotesize\color{mygray}Tuesday, 7 July 2026}\\[1pt]
{\large\bfseries N.3}\\[2pt]
{\small\textit{Ernests Ā\enspace(LV)}}\\[4pt]
{\footnotesize\color{mygray}%
  \textbf{Performance:} 18:00\,–\,open end\enspace·\enspace
  \textbf{Group visit:} 19:00}

\vspace{4pt}

{\small
\textit{The binding conviction in the existence of two opposites always makes one
wander in confusion.}

\medskip
Ernests will arrive with a rope, initially wrapping it around one of the gallery's
columns. He will then tie the end to his leg and walk around the column, attempting
to distance himself by making wide circles both within the gallery space and beyond
its walls. This action will continue until physical exhaustion or until he is stopped
by external factors. If he manages to break free from the column, he may attempt to
wrap the rope around it once again. Should Ernests need a momentary rest, he will do
so without fully interrupting the piece. He will remain silent throughout.
}

\vspace{3pt}
\sectiontitle{About the artist}

{\footnotesize\setlength{\parskip}{0pt}
In his practice, Ernests Ā devotes himself to playful exploration of the mundane and
undefinable aspects of reality within unrestricted forms of expression.
}

\vfill
\noindent{\footnotesize\color{mygray}%
  \href{https://raa.space/item/N_3}{\texttt{raa.space/item/N\_3}}\hfill\thepage}

%% ══════════════════════════════════════════════════════════
%%  PAGES ix–x: WORDS TO CHANGE THE WORLD — Ilze Aulmane  (Sun 5.07)
%% ══════════════════════════════════════════════════════════

\newpage
\newgeometry{top=0cm,bottom=0cm,left=0cm,right=0cm}
\fullbleedimage{\assets/05_07_words_to_change_the_world.png}{\thepage}

\newpage
\newgeometry{top=0.4cm,bottom=0.4cm,left=0.5cm,right=0.5cm}
\descpagebg

{\footnotesize\color{mygray}Sunday, 5 July 2026}\\[1pt]
{\large\bfseries Words to Change the World}\\[2pt]
{\small\textit{Ilze Aulmane\enspace(LV)}}\\[4pt]
{\footnotesize\color{mygray}%
  \textbf{Performance:} 19:00\,–\,20:30\enspace·\enspace
  \textbf{Group visit:} 19:30}

\vspace{4pt}

{\small
\textit{30 minutes.}

\medskip
The three words of the title are an invitation to question, challenge, and reflect.
Do we need, want, or have the power to change the world? And are there words capable
of doing so? What do we mean when we speak of ``changing the world'' — politics,
ecosystems, social environments, or our own perceptions and ways of living?

During the performance, Ilze Aulmane will create a space for this invitation through
breath, sound, letters, syllables, unfamiliar words, and performative objects.
}

\vspace{3pt}
\sectiontitle{About the artist}

{\footnotesize\setlength{\parskip}{0pt}
Ilze Aulmane is a cross-disciplinary artist working across installation, photography,
painting, sound, performance and moving image. She has developed her own experimental
vocal technique, drawing on ideas related to the origins of scat singing and the
enduring human impulse to create songs spontaneously in everyday life.

In her practice, Aulmane employs the logic of perpetual adolescence, where fragmented
and layered images and objects — both visual and acoustic — shifting bodily presences,
and movements of perception form restless spatial compositions in which the personal
and the universal continuously overlap.

Ilze Aulmane received both BA and MA degrees from the Art Academy of Latvia and
graduated from the Royal Academy Schools, London, in 2024. Alongside her academic
pursuits, she has maintained a longstanding interest in sound and has studied the
fundamentals of both classical and jazz vocals.
}

\vfill
\noindent{\footnotesize\color{mygray}%
  \href{https://raa.space/item/words-to-change-the-world}%
  {\texttt{raa.space/item/words-to-change-the-world}}\hfill\thepage}

%% ══════════════════════════════════════════════════════════
%%  PAGES xi–xii: GLUTTONY — Kristiāna Šuksta & Sandris Muriņš  (Wed 8.07)
%% ══════════════════════════════════════════════════════════

\newpage
\newgeometry{top=0cm,bottom=0cm,left=0cm,right=0cm}
\fullbleedimage{\assets/08_07_gluttony.png}{\thepage}

\newpage
\newgeometry{top=0.4cm,bottom=0.4cm,left=0.5cm,right=0.5cm}
\descpagebg

{\footnotesize\color{mygray}Wednesday, 8 July 2026}\\[1pt]
{\large\bfseries Gluttony}\\[2pt]
{\small\textit{Kristiāna Šuksta \& Sandris Muriņš\enspace(LV)}}\\[4pt]
{\footnotesize\color{mygray}%
  \textbf{Performance:} 18:00\,–\,22:00\enspace·\enspace
  \textbf{Group visit:} 20:00}

\vspace{4pt}

{\small
\textit{Gluttony} is an experimental aromagraphy performance exploring meat
consumption through the combination of poetry and scent. Built from twelve words,
each paired with its own distinct aroma, the work continuously varies the sequence,
tempo, and pauses between reading and smelling. Through both congruent and incongruent
relationships between language and scent, the performance reveals how olfactory
experience can reinforce, contradict, or transform the interpretation of a text —
exploring desire, memory, and the sensory dimensions of consumption.

By constantly reshuffling the order and timing of words and scents, the performance
disrupts linear chronology, creating a form of shuffled time where meaning emerges
through ever-changing configurations rather than a fixed sequence. Each iteration
reorders memory and anticipation, allowing past associations and present sensations
to coexist in new and unpredictable ways.
}

\vspace{3pt}
\sectiontitle{About the artists}

{\footnotesize\setlength{\parskip}{0pt}
\textbf{Sandris Muriņš} is a Latvian artist, inventor, and researcher working at
the intersection of scent, technology, and contemporary art. He is the inventor of
the Scent Camera and the founder and director of the Scent + Art Festival, actively
promoting aromagraphy as an emerging artistic discipline.

\vspace{3pt}
\textbf{Kristiāna Šuksta} is a Latvian poet, performance artist, and technical
producer of the performance festival Starptelpa. Her practice combines poetry and
live performance to explore personal experience and human relationships through wit
and humour, inviting viewers to reflect while maintaining a sense of lightness and
critical awareness.
}

\vfill
\noindent{\footnotesize\color{mygray}%
  \href{https://raa.space/item/gluttony}{\texttt{raa.space/item/gluttony}}\hfill\thepage}

%% ══════════════════════════════════════════════════════════
%%  PAGES xiii–xiv: TIME SCORES — Vanessa Brzić  (Fri 10.07)
%% ══════════════════════════════════════════════════════════

\newpage
\newgeometry{top=0cm,bottom=0cm,left=0cm,right=0cm}
\fullbleedimage{\assets/10_07_timescores.png}{\thepage}

\newpage
\newgeometry{top=0.4cm,bottom=0.4cm,left=0.5cm,right=0.5cm}
\descpagebg

{\footnotesize\color{mygray}Friday, 10 July 2026}\\[1pt]
{\large\bfseries Time Scores}\\[2pt]
{\small\textit{Vanessa Brzić\enspace(HR)}}\\[4pt]
{\footnotesize\color{mygray}%
  \textbf{Performance:} 18:00\,–\,21:00\enspace·\enspace
  \textbf{Group visit:} 19:00}

\vspace{4pt}

{\small
The performance unfolds inside a room that changes over time. At the beginning, the
space contains a body, a set of objects, and a series of event-score cards: short,
open-ended instructions that invite actions such as cutting, tying, spilling,
recording, and producing sound. Audience members draw and perform these scores,
gradually transforming the room through gestures, traces, sounds, relations, and
disruptions.

As the work progresses, the same kinds of actions return, but never under the same
conditions. An instruction such as ``cut'' or ``spill'' shifts from pure potentiality
into an action whose meaning depends on the relations already formed in space and
time. Each action inherits the consequences of previous actions while opening new
possibilities for what can happen next.

Inspired by Fluxus event scores, the piece treats the room as a metaphor for social
and historical time: a space shaped by accumulation, instability, memory, and
adjustment.
}

\vspace{3pt}
\sectiontitle{About the artist}

{\footnotesize\setlength{\parskip}{0pt}
Vanessa Brzić works across theoretical physics, experimental theatre, photography,
and performance. Originally from Zagreb, she is currently based in Paris, where she
is pursuing a PhD in Foundations of Quantum Physics and Quantum Information Theory
at Sorbonne Université. For ten years she was part of the ZKM ensemble in Zagreb,
exploring ways of deconstructing conventional performance forms. In 2024, through
an art-science residency, she began working with conceptual artist Ben Bogart.
Brzić is increasingly drawn toward conceptual and performance art, inspired by Fluxus
and process philosophy.
}

\vfill
\noindent{\footnotesize\color{mygray}%
  \href{https://raa.space/item/timescores}{\texttt{raa.space/item/timescores}}\hfill\thepage}

%% ══════════════════════════════════════════════════════════
%%  PAGES xv–xvi: LAIKS TELPĀ — Ilze Mazpane & Jana Zariņa  (Sat 11.07)
%% ══════════════════════════════════════════════════════════

\newpage
\newgeometry{top=0cm,bottom=0cm,left=0cm,right=0cm}
\fullbleedimage{\assets/11_07_laiks_telpa.png}{\thepage}

\newpage
\newgeometry{top=0.4cm,bottom=0.4cm,left=0.5cm,right=0.5cm}
\descpagebg

{\footnotesize\color{mygray}Saturday, 11 July 2026}\\[1pt]
{\large\bfseries Laiks Telpā}\\[2pt]
{\small\textit{Ilze Mazpane \& Jana Zariņa\enspace(LV)}}\\[4pt]
{\footnotesize\color{mygray}%
  \textbf{Performance:} 16:00\,–\,open end\enspace·\enspace
  \textbf{Intervention:} 18:00}

\vspace{4pt}

{\footnotesize\setlength{\parskip}{0pt}
A durational performance that transforms the existing RAA space, using only the
objects already present, into a musical, physical, and spatial experience of
subjective time. The work explores how time expands, contracts, rushes, comes to a
standstill, and disappears.

The performance begins before the performers physically enter the space and continues
with the live presence of Ilze (action) and Jana (sound) together with everyone
present.

Throughout the performance, entry and exit are possible at any time. Participation
is entirely optional and may take any form — or none at all. The main invitation is
to cultivate sensitivity to the space and attentiveness to time. Everything else
will unfold in the moment.
}

\vspace{3pt}
\sectiontitle{About the artists}

{\footnotesize\setlength{\parskip}{0pt}
\textbf{Ilze Mazpane} is a performance artist from Rīga with a background in
philosophy, visual art, experimental theatre, and creative writing. She creates
site-specific durational works marked by poetic sensibility, interest in structure,
and a sensitive relationship with space. In her practice she seeks freedom from the
tyranny of clock-time, exploring subjective time, \textit{kairos}, and the concept
of flow.

\vspace{3pt}
\textbf{Jana Zariņa} is primarily a historical keyboard and sound artist whose work
spans early music, improvisation, electro-acoustic performance, and interdisciplinary
arts. Notable projects include the electro-acoustic duo \foreignlanguage{russian}{maяk}, the
contemporary Latvian opera \textit{Ticks} (2024, Rīga Zoo), French Baroque opera
productions in Vilnius and Lithuania (2025), and the solo album
\textit{svobodulitce}. She has performed with the Latvian National Symphony Orchestra
and collaborated internationally with dancers, visual artists, and musicians.
}

\vfill
\noindent{\footnotesize\color{mygray}%
  \href{https://raa.space/item/laiks_telpa}{\texttt{raa.space/item/laiks\_telpa}}\hfill\thepage}

%% ══════════════════════════════════════════════════════════
%%  BACK COVER — full-bleed program.png (unnumbered)
%% ══════════════════════════════════════════════════════════
\newpage
\newgeometry{top=0cm,bottom=0cm,left=0cm,right=0cm}
\fullbleedimage{\assets/program.png}{}

\end{document}
